\chapter{Chikungunya Virus}
\index{Chikungunya Virus}

\section*{Definition and Overview}
The chikungunya virus is the infectious agent that causes chikungunya fever, an illness spread to humans by the bite of infected \textit{Aedes} mosquitoes. The virus is an alphavirus, first identified in Tanzania in the 1950s, and it causes a syndrome of sudden fever, severe joint and muscle pain, and often a rash. The virus does not spread directly from person to person, but an infected person can transmit it to local mosquitoes during the viraemic phase, and these mosquitoes can then pass the infection on to others. Understanding the virus's transmission cycle is central to both clinical care and outbreak control.

\section*{Epidemiology}
The chikungunya virus is found throughout tropical and subtropical regions of Africa, Asia, the Pacific, and the Americas. Large epidemics have occurred in recent decades, including in the Indian Ocean islands, Southeast Asia, and the Caribbean and Latin America. Outbreaks spread rapidly in urban areas because the vectors \textit{Aedes aegypti} and \textit{Aedes albopictus} breed in containers of standing water close to homes. Travellers frequently introduce the virus to new regions, and imported cases are reported each year in countries outside the tropics. In suitable climates, imported cases can seed local outbreaks, which makes the virus a continuing global concern.

\section*{Aetiology and Pathophysiology}
\medskip
\textbf{Transmission and infection.} The virus enters the human body when an infected mosquito takes a blood meal:
\begin{itemize}
    \item \textbf{Primary replication:} the virus first multiplies in the skin and lymphatic tissue at the site of the bite
    \item \textbf{Bloodstream spread:} it then disseminates through the bloodstream to reach the joints, muscles, and skin
    \item \textbf{Cellular infection and inflammation:} once inside cells, the virus triggers an intense inflammatory response that produces fever and the characteristic joint pain and swelling
    \item \textbf{Viral clearance:} the immune system normally clears the virus within one to two weeks
    \item \textbf{Chronic symptoms:} in some cases, viral fragments or a persisting inflammatory process remain in the joints, which is thought to explain chronic pain long after the acute illness
\end{itemize}

\section*{Clinical Features}
\begin{itemize}
    \item \textbf{Incubation:} symptoms begin abruptly, usually two to six days after an infected mosquito bite
    \item \textbf{Fever:} sudden onset of high fever
    \item \textbf{Joint pain:} severe, most often affecting the wrists, fingers, ankles, and knees
    \item \textbf{Joint swelling and stiffness:} commonly on both sides of the body
    \item Muscle aches, headache, and fatigue
    \item A flat red rash over the trunk and limbs in many patients
    \item Nausea or vomiting in some cases
    \item Symptoms can develop within hours and are often incapacitating at the peak
\end{itemize}

\section*{Differential Diagnosis}
\begin{itemize}
    \item Dengue virus infection -- fever, headache, and muscle aches, with a higher risk of bleeding problems
    \item Zika virus infection -- generally milder, with a prominent rash
    \item Ross River and Barmah Forest virus infections, seen in many countries
    \item Malaria, which must always be considered in returning travellers
    \item Influenza and other viral illnesses with fever and body aches
    \item Rubella or parvovirus infection presenting with a febrile rash
\end{itemize}

\section*{When to Seek Medical Care}
Anyone with fever and joint pain after travel to an affected region should seek medical assessment, particularly if symptoms are severe or accompanied by other worrying features.

\medskip
\textbf{Contact your local emergency services for:}
\begin{itemize}
    \item High fever that does not settle, confusion, or a severe headache
    \item Difficulty breathing or chest pain
    \item Vomiting, poor fluid intake, or signs of dehydration
    \item Bleeding from the gums or nose, or easy bruising
    \item Infection in an infant, an older person, or a person with chronic disease
\end{itemize}

\section*{Diagnosis and Investigations}
\begin{itemize}
    \item \textbf{History and travel history:} exposure to mosquitoes and time spent in affected regions
    \item \textbf{Blood PCR:} detects the virus itself during the first week of illness
    \item \textbf{Serology:} antibody tests (IgM and IgG) become positive after about a week and confirm recent or past infection
    \item \textbf{Full blood count:} may show low white cells and platelets
    \item \textbf{Liver and inflammatory markers:} often mildly abnormal
    \item \textbf{Tests for dengue and malaria:} needed to exclude other infections with similar presentations
    \item \textbf{Public health notification:} laboratory-confirmed cases are reported so that local transmission can be detected and controlled
\end{itemize}

\section*{Management}
\begin{itemize}
    \item \textbf{Supportive care:} rest, fluids, and paracetamol for fever in the recommended dose
    \item \textbf{Joint pain relief:} non-steroidal anti-inflammatory medicines such as ibuprofen after dengue has been excluded
    \item \textbf{Physiotherapy:} gentle mobilisation and exercise to help joints recover and prevent stiffness
    \item \textbf{Follow-up:} review for persistent joint or fatigue symptoms
    \item \textbf{No specific antiviral:} treatment targets symptoms while the immune system clears the virus
    \item \textbf{Mosquito precautions during illness:} avoiding mosquito bites prevents transmission back to local mosquitoes
\end{itemize}

\section*{Complications}
\begin{itemize}
    \item Chronic, sometimes disabling joint pain lasting weeks to years
    \item Long-term fatigue and reduced mobility, particularly in older patients
    \item Rare heart and eye involvement during the acute illness
    \item Neurological complications, including encephalitis, are uncommon but reported
    \item Severe infection in newborns exposed around the time of birth
\end{itemize}

\section*{Prognosis and Outcomes}
The acute illness usually settles within one to two weeks, and most people recover completely. A significant proportion of patients, particularly older adults, experience joint pain that continues for months, and some have symptoms lasting a year or more, which can interfere with work and daily activities. Death is rare and is most likely in the very young, the very old, and people with underlying disease. Because infection produces long-lasting immunity, repeat attacks are uncommon, and most patients can expect eventual full recovery.

\section*{Prevention and Self-Care}
\begin{itemize}
    \item Apply insect repellent to exposed skin during the day, when the vector mosquitoes are most active
    \item Cover up with long sleeves and long trousers in affected areas
    \item Sleep under mosquito nets and use screened windows and doors
    \item Empty, cover, or treat standing water around the home to remove breeding sites
    \item Travellers should plan bite-avoidance measures before visiting outbreak regions
    \item Ask a doctor about vaccination against chikungunya if travelling to high-risk areas
    \item People with an active infection should avoid mosquito bites until fully recovered
\end{itemize}

\section*{Sources and Further Reading}
The following reputable sources provide further information on the chikungunya virus:
\begin{itemize}
    \item World Health Organization. \textit{Chikungunya fact sheets and technical guidance.} https://www.who.int/
    \item Centers for Disease Control and Prevention. \textit{Chikungunya virus resources.} https://www.cdc.gov/
    \item NHS. \textit{Health A--Z: travel health and infectious disease information.} https://www.nhs.uk/conditions/
    \item MSD Manual. \textit{Professional edition: arbovirus infections.} https://www.msdmanuals.com/
    \item MedlinePlus. \textit{Health topics on infectious diseases.} https://medlineplus.gov/
\end{itemize}
